% ----------------------------------------------------------------------
% Apêndice A: o produto técnico
% Apêndice é material elaborado pelo próprio autor (NBR 14724:2011, 4.2.3.3).
% ----------------------------------------------------------------------
\chapter{Produto técnico: guia de política de desenvolvimento de coleções}
\label{ap:produto}

\instrucao{Coloque aqui o produto completo, no formato em que será entregue à instituição. Se ele tiver identidade visual própria, mantenha a página de rosto do produto e a ficha técnica. Produtos em outro suporte (vídeo, sistema, site) entram com descrição, endereço de acesso e capturas de tela.}

\begin{center}
  \fcolorbox{azulbib}{azulbib!6}{\begin{minipage}{0.86\textwidth}
    \centering\vspace{1em}
    {\large\bfseries\color{azulbib} Guia de política de desenvolvimento de coleções\par}
    \medskip
    {Bibliotecas escolares da rede pública municipal\par}
    \medskip
    {\small Nome do autor e Nome da orientadora\par Cidade, Ano\par}
    \vspace{1em}
  \end{minipage}}
\end{center}

\section*{Ficha técnica}
\noindent\begin{tabularx}{\textwidth}{@{}lX@{}}
  Tipo de produto & Manual técnico com formulários \\
  Público & Equipes de bibliotecas escolares \\
  Registro & ISBN ou repositório institucional, quando houver \\
  Licença & Indique a licença de uso, por exemplo CC BY 4.0 \\
\end{tabularx}

\section*{1 Objetivos da política}
\instrucao{Texto-modelo que a escola completa: missão da biblioteca, público e níveis de ensino atendidos.}
A biblioteca da Escola \underline{\hspace{4cm}} atende estudantes do ensino fundamental, professores e comunidade, e tem por missão apoiar o currículo e incentivar a leitura.

\section*{2 Critérios de seleção}
\begin{enumerate}[label=\alph*)]
  \item adequação à faixa etária e ao currículo;
  \item qualidade do conteúdo e da edição;
  \item atualidade, nas áreas em que a informação envelhece depressa;
  \item demanda registrada no livro de sugestões.
\end{enumerate}

\section*{3 Doações}
Toda doação é recebida mediante termo, e a biblioteca decide se incorpora, repassa ou descarta o material segundo os critérios da seção 2.

\section*{4 Avaliação e desbaste}
O acervo é avaliado a cada dois anos. Materiais desatualizados, danificados ou sem empréstimo há cinco anos são remanejados ou descartados, com registro em formulário próprio.

\section*{Formulário 1: termo de doação}
\noindent Doador(a): \hrulefill\par\smallskip
\noindent Quantidade de itens: \underline{\hspace{2cm}}\quad Data: \underline{\hspace{3cm}}\par\smallskip
\noindent Declaro estar ciente de que a biblioteca decidirá sobre a destinação dos itens.\par\bigskip
\noindent\hfill\rule{6cm}{0.4pt}\par
\noindent\hfill Assinatura do doador
